\documentclass[11pt,a4paper]{article}

%================================================
% ENCODAGE ET LANGUE
%================================================

\usepackage[utf8]{inputenc}
\usepackage[T1]{fontenc}
\usepackage[french]{babel}

%================================================
% MISE EN PAGE
%================================================

\usepackage{geometry}
\geometry{margin=1.55cm}

%================================================
% PACKAGE GTOLI
%================================================

\usepackage{../styles/gtoli}

%================================================
% INFORMATIONS DU DOCUMENT
%================================================

\title{\textbf{Titre du chapitre}}
\author{GToli}
\date{\today}

%================================================
% DOCUMENT
%================================================

\begin{document}

\maketitle

%================================================
% INFORMATIONS GÉNÉRALES
%================================================

\begin{cadreobjectif}
\begin{itemize}
    \item Comprendre les notions essentielles du chapitre ;
    \item maîtriser les méthodes de résolution ;
    \item développer le raisonnement ;
    \item appliquer les concepts dans des exercices.
\end{itemize}
\end{cadreobjectif}

%================================================
\section{Notion concernée}
%================================================

\begin{cadrebleu}{Introduction}
Cette section permet d’introduire la notion étudiée et de contextualiser le chapitre.
\end{cadrebleu}

%================================================
\section{Rappel théorique}
%================================================

\begin{cadretheorie}
Cette section contient :
\begin{itemize}
    \item définitions ;
    \item propriétés ;
    \item théorèmes ;
    \item formules ;
    \item vocabulaire essentiel.
\end{itemize}
\end{cadretheorie}

%================================================
\section{Méthode / protocole}
%================================================

\begin{cadremethode}
Cette section présente :
\begin{enumerate}
    \item les étapes ;
    \item la démarche ;
    \item le raisonnement ;
    \item les points d’attention ;
    \item les erreurs fréquentes.
\end{enumerate}
\end{cadremethode}

%================================================
\section{Exemple guidé}
%================================================

\begin{cadreexemple}
Cette section contient un exemple entièrement expliqué étape par étape.
\end{cadreexemple}

%================================================
\section{Exercice d’application}
%================================================

\begin{cadreenonce}
Cette section contient un exercice ou une activité.
\end{cadreenonce}

%================================================
\section{Aide / piste}
%================================================

\begin{cadreaide}
Cette section contient :
\begin{itemize}
    \item indices ;
    \item rappels ;
    \item aides méthodologiques ;
    \item pistes de réflexion.
\end{itemize}
\end{cadreaide}

%================================================
\section{Solution détaillée}
%================================================

\begin{cadresolution}
Cette section contient :
\begin{itemize}
    \item les calculs ;
    \item les justifications ;
    \item les explications ;
    \item les interprétations.
\end{itemize}
\end{cadresolution}

%================================================
\section{À retenir}
%================================================

\begin{cadresynthese}
Cette section synthétise :
\begin{itemize}
    \item les idées importantes ;
    \item les formules essentielles ;
    \item les méthodes à connaître ;
    \item les points de vigilance.
\end{itemize}
\end{cadresynthese}

%================================================
\section{Exercices supplémentaires}
%================================================

\begin{cadreactivite}
\begin{enumerate}
    \item Exercice 1
    \item Exercice 2
    \item Exercice 3
\end{enumerate}
\end{cadreactivite}

%================================================
\section{Auto-évaluation}
%================================================

\begin{cadregris}{Questions de réflexion}
\begin{itemize}
    \item Suis-je capable d’expliquer la méthode ?
    \item Suis-je capable de refaire l’exercice seul ?
    \item Ai-je compris les erreurs fréquentes ?
    \item Puis-je résoudre une variante ?
\end{itemize}
\end{cadregris}

\end{document}